\documentclass[14pt, a4paper]{extarticle}
\usepackage[T2A]{fontenc}
\usepackage[utf8]{inputenc}
\usepackage[english, russian]{babel}
\usepackage[left=30mm, right=10mm, top=20mm, bottom=25mm]{geometry}
\usepackage{amsmath, amssymb, amsfonts}
\usepackage{graphicx}
\usepackage{float}
\usepackage{hyperref}
\usepackage{longtable}
\usepackage{booktabs}
\usepackage{multirow}
\usepackage{enumitem}
\usepackage{array}
\usepackage{setspace}

\usepackage{misccorr}
\usepackage{indentfirst}

\setstretch{1.2}

% Заголовки section по центру
\makeatletter
\renewcommand\section{\@startsection {section}{1}{\z@}%
{-3.5ex \@plus -1ex \@minus -.2ex}%
{2.3ex \@plus.2ex}%
{\centering\normalfont\Large\bfseries}}
\makeatother

% Точки в списке литературы
\makeatletter

\makeatother

% Настройка hyperref
\hypersetup{
colorlinks=true,
linkcolor=black,
citecolor=black,
urlcolor=blue
}

\begin{document}

\thispagestyle{empty}
{\scriptsize
\begin{center}
ФЕДЕРАЛЬНОЕ ГОСУДАРСТВЕННОЕ ОБРАЗОВАТЕЛЬНОЕ БЮДЖЕТНОЕ УЧРЕЖДЕНИЕ\\ 
ВЫСШЕГО ОБРАЗОВАНИЯ    
\end{center}
}
{\bf
\begin{center}
ФИНАНСОВЫЙ УНИВЕРСИТЕТ\\ 
ПРИ ПРАВИТЕЛЬСТВЕ РОССИЙСКОЙ ФЕДЕРАЦИИ\\
Липецкий филиал
\end{center}
}
\hrule
\bigskip

{\small 
\begin{center}
\textbf{
Кафедра <<Учет и информационные технологии в бизнесе>>}
\end{center}
}
\bigskip

Направление: Прикладная математика и информатика

\medskip
Профиль: Анализ данных и принятие решений в экономике и финансах

\vskip 2cm

\begin{center}
\textbf{КУРСОВОЙ ПРОЕКТ}
\end{center}

По дисциплине: <<Машинное обучение>>

\bigskip
ТЕМА: Проектирование и реализация клиентской части веб-приложения 
FITGUIDE с поддержкой интерактивных компонентов машинного обучения

\vskip 3cm

\hskip .3\textwidth Преподаватель: Калитвин Владимир Анатольевич

\hskip .3\textwidth Студент: Королёв Артём Игоревич

\hskip .3\textwidth Учетный номер \textnumero 100.16/240080

\hskip .3\textwidth Курс второй

\vskip 6cm

\begin{center}
Липецк 2026
\end{center}

\newpage
\tableofcontents
\newpage

\addcontentsline{toc}{section}{ВВЕДЕНИЕ}
\section*{ВВЕДЕНИЕ}

\noindent Курсовой проект по машинному обучению - часть программы на втором курсе направления «Прикладная информатика». Многие одногруппники сначала думали: ML - это запредельно сложно, доступно только учёным с многолетним стажем. Но уже после пары лекций и практики выяснилось: современные технологии дают готовые нейросети через API. Значит, основная задача сделать полезное и удобное приложение, а не писать нейросеть с нуля. Это значительно упрощает задачу и позволяет сосредоточиться на самом главном — создании полезного и удобного приложения.

\medskip
\noindent Долго думалось над темой. Хотелось сделать не просто работу «для галочки», а реально полезную вещь. Многие ходят в зал, и постоянно мучаются вопросом: где взять нормальную программу тренировок и питания? Так родилась идея фитнес-помощника. В интернете существует огромное количество статей и видео на тему фитнеса, но почти все они дают общие рекомендации, не учитывающие индивидуальные особенности человека: вес, рост, тип телосложения, уровень подготовки и, самое главное, конкретную цель. Именно эту проблему решит данный проект.

\medskip
\noindent Так появилась идея создания веб-приложения FITGUIDE. Пользователь вводит свои физические параметры и желаемую цель (например, «набрать мышечную массу» или «похудеть к лету»), а нейросеть генерирует персональный план тренировок и питания. Ключевая особенность проекта заключается в том, что разрабатывается не просто «прикрученная» нейросеть к какому-то шаблону, а полноценная клиентская часть - интерфейс, с которым пользователь будет взаимодействовать. Именно клиентская часть определяет, насколько приложение будет удобным, понятным и привлекательным.

\medskip
\noindent До этого курсового проекта не было опыта в веб-разработке, поэтому это был абсолютно новый опыт. Предстояло разобраться с тем, как устроены веб-приложения, как отправлять запросы к удалённым серверам, как обрабатывать ответы нейросети и самое главное, как сделать всё это красиво и удобно для пользователя.

\subsection*{Актуальность темы}

\noindent Фитнес сейчас реально в моде. Люди всё больше втягиваются в спорт, следят за здоровьем и внешностью. Но цифры жестокие: больше половины новичков бросают тренировки в первые три месяца. Почему? Одна из главных причин - отсутствие персонального плана и понятной системы.

\medskip
\noindent Но не каждый может себе позволить тренера или диетолога. В крупных городах час тренировки - от двух тысяч рублей. Консультация диетолога - от пяти до десяти тысяч. Дорого. Готовые программы из интернета часто не учитывают индивидуальные особенности, а порой и вовсе могут навредить.

\medskip
\noindent С развитием технологий машинного обучения появилась возможность создавать персонализированные рекомендации автоматически. Современные языковые модели способны генерировать связные и осмысленные предложения, учитывающие множество входных параметров. При этом важно, чтобы интерфейс такого приложения был максимально простым и понятным, чтобы им мог пользоваться любой человек, даже не имеющий специальной подготовки.

\medskip
\noindent Таким образом, разработка удобной клиентской части для ML-приложений является актуальной задачей, имеющей как теоретическое, так и практическое значение. Данный проект может служить примером того, как современные технологии делают фитнес более доступным.

\subsection*{Цель работы}

\noindent Целью данной курсовой работы является проектирование и программная реализация клиентской части веб-приложения FITGUIDE, которая обеспечивает решение следующих задач:
\begin{itemize}
\item сбор физических параметров пользователя с использованием различных типов полей ввода (числовые значения, выпадающие списки, ползунки, множественный выбор);
\item ввод цели в свободной текстовой форме с поддержкой подсказок и примеров для заполнения;
\item отправка данных на сервер и получение ответа от модели машинного обучения через API с отображением индикатора загрузки;
\item наглядное представление полученных результатов в эстетичной форме с опорой на современные принципы интерфейсного дизайна;
\item сохранение сформированного плана в текстовый файл для последующего использования пользователем;
\item показ случайного совета дня, который обновляется при каждой перезагрузке страницы.
\end{itemize}

\subsection*{Задачи работы}

\noindent Для достижения поставленной цели необходимо решить следующие задачи:

\begin{enumerate}
\item Изучить современные подходы к созданию пользовательских интерфейсов для приложений с использованием машинного обучения. Рассмотреть особенности, отличающие ML-приложения от традиционных веб-сервисов.

\item Провести обзор и сравнение инструментов, доступных для быстрой разработки клиентской части на языке Python. Выделить преимущества и недостатки каждого из рассмотренных решений.

\item Обосновать выбор технологического стека (фреймворка, API, модели) для реализации проекта с учётом требований к бесплатности, простоте использования и качеству генерации на русском языке.

\item Разработать структуру и макеты интерфейса приложения, продумать сценарии использования для разных категорий пользователей.

\item Реализовать клиентскую часть с использованием выбранных технологий, уделяя особое внимание управлению состоянием приложения и обработке ошибок.

\item Выполнить кастомизацию внешнего вида с применением каскадных таблиц стилей (CSS) для придания приложению современного и привлекательного облика.

\item Провести тестирование разработанного интерфейса, выявить и исправить ошибки, оптимизировать производительность.

\item Оформить отчёт о проделанной работе в соответствии с требованиями кафедры и ГОСТ.

\subsection*{Объект и предмет исследования}

\noindent \textbf{Объектом исследования} являются процессы взаимодействия пользователя с веб-приложениями на основе машинного обучения. В рамках данной работы изучаются способы организации ввода данных, отправки запросов к удалённым API и отображения результатов.

\medskip
\noindent \textbf{Предметом исследования} выступает клиентская часть веб-приложения FITGUIDE, включающая компоненты ввода данных, отправки запросов к API нейросетевой модели и отображения полученных результатов. Особое внимание уделяется вопросам удобства использования и визуальному оформлению.

\subsection*{Методы исследования}

\noindent При выполнении работы использовались следующие методы:
\begin{itemize}
\item анализ научно-технической литературы и документации по выбранным технологиям;
\item сравнительный анализ инструментов разработки;
\item метод прототипирования при создании интерфейса;
\item экспериментальное тестирование разработанного приложения;
\item метод наблюдения за взаимодействием пользователя с интерфейсом.
\end{itemize}
\vskip 5cm
\subsection*{Структура работы}

\noindent Курсовая работа состоит из введения, пяти глав, заключения, списка литературы (одиннадцать источников) и приложений. Общий объём работы составляет 38 страниц машинописного текста.

\medskip
\noindent В первой главе рассматриваются теоретические основы разработки клиентской части ML-приложений, включая определение клиентской части, особенности интерфейсов для машинного обучения и проблемы хранения состояния.

\medskip
\noindent Во второй главе приводится обзор существующих решений для фитнес-планирования и сравнительный анализ инструментов разработки, обосновывается выбор технологического стека.

\medskip
\noindent Третья глава посвящена проектированию интерфейса FITGUIDE: описанию функциональных требований, разработке макетов, цветовой схемы и сценариев использования.

\medskip
\noindent В четвёртой главе описывается процесс реализации клиентской части, включая вопросы управления состоянием, подключения к API, кастомизации дизайна и обработки ошибок.

\medskip
\noindent Пятая глава содержит результаты тестирования, выявленные проблемы и их решения, а также экономическую оценку проекта.

\newpage
\section{ТЕОРЕТИЧЕСКИЕ ОСНОВЫ РАЗРАБОТКИ КЛИЕНТСКОЙ ЧАСТИ ML-ПРИЛОЖЕНИЙ}

\subsection{Понятие клиентской части веб-приложения}

\noindent Для того чтобы лучше понимать, о чём идёт речь в данной работе, необходимо сначала разобраться с базовыми понятиями. Любое современное веб-приложение можно разделить на две большие части: клиентскую (фронтенд) и серверную (бэкенд). Это разделение является фундаментальным для понимания архитектуры веб-приложений.

\medskip
\noindent \textbf{Клиентская часть} - это всё, что видит и с чем взаимодействует пользователь, открывая веб-страницу в браузере. Клиентская часть включает в себя все визуальные и интерактивные элементы. К клиентской части относятся:

\begin{itemize}
\item текстовые поля, кнопки, переключатели, ползунки — все элементы управления, с помощью которых пользователь вводит данные и управляет приложением;
\item изображения, анимации, иконки - все визуальные элементы, которые делают интерфейс привлекательным и понятным;
\item расположение элементов на странице (вёрстка) - то, как элементы организованы в пространстве, какие между ними отступы и порядок;
\item цвета, шрифты, тени и другие визуальные эффекты - стилизация, которая создаёт настроение и помогает пользователю ориентироваться;
\item поведение интерфейса при наведении мыши, нажатии клавиш - интерактивные реакции, которые дают пользователю обратную связь.
\end{itemize}
\vskip 5cm
\noindent \textbf{Серверная часть} - это то, что работает на удалённом компьютере (сервере) и обычно не видно пользователю. Серверная часть выполняет «тяжёлую» работу, которая требует вычислительных ресурсов или доступа к базам данных. К серверной части относятся:

\begin{itemize}
\item базы данных, в которых хранится информация о пользователях, их действиях, сохранённых настройках;
\item сложные вычисления и обработка данных, которые требуют большой
вычислительной мощности;
\item модели машинного обучения, которые занимают много памяти и требуют специального оборудования (например, видеокарт) для работы;
\item бизнес-логика приложения - правила, по которым приложение принимает решения;
\item обработка запросов от клиентской части - сервер слушает запросы от клиентов и отправляет им ответы.
\end{itemize}

\noindent В традиционной архитектуре, которую часто называют «клиент-серверной», клиентская часть отправляет запросы на сервер, сервер их обрабатывает и возвращает результат обратно клиенту. Клиентская часть только отображает этот результат и не занимается сложными вычислениями. Такая архитектура надёжна и масштабируема, но требует настройки и поддержки сервера.

\medskip
\noindent Однако в данном проекте используется несколько иная схема. Был выбран путь разработки собственного сервера с базой данных и логикой обработки, так как это потребовало бы значительно больше времени и ресурсов. Вместо этого использовано готовое решение — Hugging Face Inference API.

\medskip
\noindent В этой архитектуре клиентская часть напрямую отправляет запросы к API нейросети, получает ответ и сама его отображает. Промежуточного собственного сервера нет. Такой подход называют «тонкий фронтенд плюс облачный API» или иногда «бессерверная архитектура».

\medskip
\noindent Данный подход имеет ряд важных преимуществ, которые оказались решающими для студенческого проекта:

\begin{itemize}
\item \textbf{Отсутствие необходимости настраивать и поддерживать свой сервер.} Это экономит время и деньги, не нужно покупать хостинг или арендовать виртуальную машину.

\item \textbf{Не нужно беспокоиться о масштабировании.} Hugging Face сам справляется с большим количеством запросов. Если вдруг приложением заинтересуются тысячи пользователей, ничего не нужно менять в коде.

\item \textbf{Разработка происходит быстрее и проще.} Сосредоточение только на интерфейсе, а вся сложная работа с нейросетью уже сделана профессионалами.

\item \textbf{Можно использовать самые современные нейросетевые модели.} Нет необходимости покупать дорогие видеокарты и устанавливать сложное программное обеспечение. Всё уже работает в облаке.
\end{itemize}

\noindent Однако есть и недостатки, о которых стоит упомянуть. Главный недостаток - зависимость от интернета. Если у пользователя пропадёт соединение, нейросеть не ответит. Также есть ограничения на количество бесплатных запросов в день. Для Hugging Face это около пятисот запросов, что для данного проекта (учебного, не коммерческого) вполне достаточно.

\subsection{Особенности интерфейсов для приложений машинного обучения}

\noindent Приложения, в основе которых лежат модели машинного обучения, имеют определённые особенности, которые необходимо учитывать при проектировании клиентской части. В ходе работы над проектом были выделены несколько ключевых моментов, которые отличают ML-приложения от традиционных веб-сервисов. Эти особенности связаны с самой природой нейросетей и их поведением.

\subsubsection{Неопределённое время ответа}

\noindent В обычных приложениях результат приходит за доли секунды - нажал, и готово. С нейросетями иначе: отклика надо ждать несколько секунд, а то и дольше. Потому что внутри модели - миллиарды математических операций, даже если сервер мощный.

\medskip
\noindent Пользователь, нажав на кнопку, не должен думать, что приложение «зависло» или «сломалось». Если ничего не происходит в течение трёх-пяти секунд, у большинства пользователей возникает желание закрыть страницу или нажать кнопку ещё раз, что может привести к повторной отправке запроса и ещё большей задержке.

\medskip
\noindent Поэтому необходимо предусмотреть визуальную индикацию процесса обработки. В данном приложении при отправке запроса появляется спиннер (крутящийся значок) и текст «Генерируем план...». Это даёт пользователю понять, что приложение работает, нужно немного подождать. Также кнопка становится неактивной, чтобы пользователь не мог нажать её повторно.

\subsubsection{Неопределённость и вариативность результата}

\noindent Нейросеть — это вероятностная модель. Она может выдать разный ответ на один и тот же запрос в разные моменты времени. Это особенность работы языковых моделей: они не хранят ответы в базе данных, а каждый раз генерируют текст заново, выбирая слова с некоторой вероятностью.

\medskip
\noindent Более того, иногда ответ может быть неполным, содержать ошибки или даже быть не совсем по теме. Например, на запрос про набор мышечной массы нейросеть может начать рассказывать про кардиотренировки или питание для похудения. Это связано с тем, что модель «не понимает» смысл, а лишь предсказывает наиболее вероятные слова.

\vskip 5cm
\noindent Поэтому важно предусмотреть возможность повторной генерации. Если пользователю не понравился план, он может слегка изменить формулировку цели и нажать кнопку снова. Также в интерфейс добавлены понятные сообщения, чтобы пользователь не пугался, если ответ покажется ему «странным».

\subsubsection{Необходимость обработки ошибок API}

\noindent API нейросети может быть недоступно по разным причинам. В ходе разработки встретились несколько типичных ситуаций:

\begin{itemize}
\item превышение лимита бесплатных запросов - Hugging Face даёт около пятисот запросов в день, после чего доступ блокируется до следующего дня;
\item проблемы с интернет-соединением на стороне пользователя или на стороне провайдера;
\item технические работы на серверах Hugging Face;
\item временная перегрузка API из-за большого количества пользователей;
\item ошибки в самом запросе (например, слишком длинный текст или недопустимые символы).
\end{itemize}

\noindent Приложение не должно падать с ошибкой, показывая пользователю страшные сообщения на английском языке, полные технических терминов вроде «ConnectionError» или «HTTP 503». Это неприемлемо с точки зрения пользовательского опыта.

\medskip
\noindent Вместо этого нужно применить подход, который в разработке называется «graceful degradation» - изящная деградация. Это означает, что приложение продолжает работать, пусть и в урезанном режиме, показывая заранее заготовленные ответы или понятные сообщения об ошибке на родном языке пользователя.

\vskip 5cm
\noindent В данном приложении реализовано следующее: если API не отвечает, блок «Совет дня» показывает заранее заготовленный текст, а при попытке сгенерировать план появляется сообщение: «Ошибка при обращении к нейросети. Попробуйте позже или измените формулировку цели». Это гораздо лучше, чем просто красная надпись с английскими буквами.

\subsubsection{Персонализация как ключевое преимущество}

\noindent Главное преимущество ML-приложений перед традиционными программами - это возможность учёта индивидуальных особенностей пользователя. Если обычная программа работает по жёсткому алгоритму «если А, то Б», то нейросеть способна учитывать десятки параметров одновременно и выдавать уникальный результат для каждого пользователя.

\medskip
\noindent Поэтому интерфейс должен позволять вводить как можно больше релевантных параметров, которые влияют на конечный результат. Чем больше информации получит нейросеть, тем более персонализированным и полезным будет ответ.

\medskip
\noindent В данном приложении пользователь может указать следующие параметры:
\begin{itemize}
\item вес в килограммах (число от 40 до 200);
\item рост в сантиметрах (число от 100 до 250);
\item пол (мужской или женский);
\item тип телосложения (эктоморф, мезоморф, эндоморф);
\item уровень опыта (новичок, средний, профи);
\item доступный инвентарь (зал, гантели, турник, без инвентаря - можно выбрать несколько вариантов).
\end{itemize}

\noindent Все эти параметры влияют на конечный результат. Например, эктоморфу (худому от природы человеку) подойдёт одна программа набора массы, а эндоморфу (склонному к полноте) - совершенно другая. Новичку нужны одни упражнения, профессионалу другие.

\subsection{Проблема хранения состояния в веб-приложениях}

\noindent При изучении выбранного фреймворка обнаружилась интересная особенность. Оказалось, что он работает не так, как обычные Python-скрипты, к которым привыкли за время учёбы.

\medskip
\noindent В обычном скрипте переменная создаётся один раз и живёт до конца работы программы. Пример: написан счётчик = 0, потом +1, потом ещё +1 - в переменной двойка. Всё как надо.

\medskip
\noindent Но в данном веб-фреймворке всё хитрее: при любом действии пользователя (нажал кнопку, подвинул ползунок) скрипт запускается заново - с первой строки до последней. Такая архитектура из-за того, как HTTP-запросы обрабатываются.

\medskip
\noindent Давайте рассмотрим простой пример, чтобы стало понятнее. Если попытаться сделать кнопку-счётчик, которая показывает, сколько раз на неё нажали, и написать код традиционным способом (объявив переменную и увеличивая её при нажатии), то счётчик всегда будет показывать единицу. Почему? Потому что при каждом нажатии скрипт запускается заново, переменная снова становится нулём, потом единицей, и пользователь видит «нажатий: 1». Счётчик не растёт.

\medskip
\noindent Эта проблема называется «потеря состояния». Переменные не сохраняются между разными «запусками» скрипта, которые происходят при каждом взаимодействии пользователя с интерфейсом.

\medskip
\noindent Для решения этой проблемы в выбранном фреймворке существует специальный механизм - хранилище сессии (session state). Это специальное место в памяти, которое сохраняет свои значения между перезапусками скрипта. Если переменная хранится в session state, она не теряется при нажатии на кнопку.

\medskip
\noindent В данном приложении session state используется для хранения цели, которую вводит пользователь. Это необходимо, потому что при нажатии на кнопку «Сформировать фитнес-план» скрипт перезапускается, и если бы не использовался session state, введённая цель терялась бы. Пользователю пришлось бы вводить её заново каждый раз - это ужасный пользовательский опыт.

\subsection{Понятие API и работа с нейросетевыми моделями}

\noindent Термин API расшифровывается как Application Programming Interface, что в переводе на русский означает «программный интерфейс приложения». Если говорить совсем просто, API - это способ, с помощью которого одна программа может обращаться к функциям другой программы. Это как розетка и вилка: у сервиса есть «розетка» (API) с определёнными контактами (методами), а клиентская программа подключает к этой розетке свою «вилку» и получает доступ к функциональности.

\medskip
\noindent В данном случае приложение на Python (клиентская часть) обращается к API нейросети, которая находится на серверах компании Hugging Face. Нейросеть не запускается на локальном компьютере - это потребовало бы слишком много ресурсов и специальных знаний. Вместо этого отправляется запрос на их сервер, они выполняют вычисления и присылают ответ.

\medskip
\noindent Процесс работы с API выглядит следующим образом:

\begin{enumerate}
\item Приложение формирует текст запроса (в терминологии нейросетей это называется «промпт») на основе параметров, которые ввёл пользователь. Промпт - это инструкция для нейросети, в которой объясняется, кем она должна себя считать, какие у неё есть данные и что от неё требуется.

\item Приложение отправляет HTTP-запрос к API по определённому адресу (URL). HTTP - это протокол передачи гипертекста, по которому работает весь интернет. Запрос содержит промпт и секретный токен (ключ доступа).

\item API принимает запрос, проверяет токен (есть ли право пользоваться сервисом) и передаёт промпт нейросети.

\item Нейросеть выполняет вычисления: она анализирует промпт, определяет контекст и начинает генерировать ответ слово за словом.

\item API упаковывает ответ в специальный формат (обычно JSON - текстовый формат для обмена данными) и отправляет обратно приложению.

\item Приложение получает ответ, извлекает из него нужный текст (убирая лишнюю техническую информацию) и отображает его пользователю в красивом оформленном блоке.
\end{enumerate}

\noindent Для работы с Hugging Face API использована официальная библиотека, которую предоставляет сама компания. Эта библиотека упрощает процесс отправки запросов и получения ответов, избавляя от необходимости писать низкоуровневый код для работы с HTTP.

\medskip
\noindent Важно отметить, что для использования API большинства современных нейросетей требуется регистрация и получение токена (ключа доступа). Токен нужно хранить в секрете, потому что с его помощью кто угодно может отправлять запросы от имени владельца. В данном проекте токен хранится в отдельном файле, который не попадает в публичный доступ.

\newpage
\section{ОБЗОР И ВЫБОР ТЕХНОЛОГИЙ ДЛЯ РЕАЛИЗАЦИИ ПРОЕКТА}

\subsection{Сравнение существующих решений для фитнес-планирования}

\noindent Прежде чем приступить к разработке собственного приложения, был проведён анализ существующих на рынке фитнес-приложений. Это помогло понять, какие функции являются важными для пользователей, а какие - второстепенными, а также выявить недостатки существующих решений, которые можно исправить в данном проекте. Анализ проводился по следующим критериям: функциональность, стоимость, качество интерфейса, возможность персонализации.

\subsubsection{MyFitnessPal}

\noindent MyFitnessPal - это одно из самых популярных в мире приложений для подсчёта калорий и планирования питания. Оно было запущено в 2005 году и на сегодняшний день имеет более двухсот миллионов зарегистрированных пользователей. Основные характеристики этого приложения:

\begin{itemize}
\item огромная база продуктов питания - более четырнадцати миллионов наименований, можно сканировать штрих-коды для быстрого добавления;
\item интеграция с фитнес-трекерами и умными часами (Apple Watch, Fitbit, Garmin и другие);
\item возможность ставить цели по снижению, поддержанию или набору веса;
\item дневник питания с визуализацией калорий и питательных веществ (белки, жиры, углеводы);
\item сообщество пользователей, где можно делиться успехами и получать поддержку.
\end{itemize}

\noindent \textbf{Недостатки:} многие расширенные функции доступны только по платной подписке (около 500 рублей в месяц или 3000 рублей в год). Нет персонализированных планов тренировок - приложение фокусируется только на питании. Интерфейс перегружен информацией, новым пользователям бывает сложно разобраться.

\subsubsection{Fitbod}

\noindent Fitbod - это приложение, которое специализируется на генерации тренировочных программ. Оно использует технологию машинного обучения для адаптации нагрузки под прогресс пользователя. Основные характеристики:

\begin{itemize}
\item программа тренировок автоматически подстраивается под результаты прошлых занятий;
\item визуальные инструкции к каждому упражнению с анимацией правильного выполнения;
\item синхронизация с Apple Health и Google Fit;
\item возможность указать доступное оборудование, и приложение подберёт упражнения только с ним;
\item статистика прогресса: сколько подходов сделано, какой вес поднят, как меняется сила.
\end{itemize}

\noindent \textbf{Недостатки:} платная подписка - около 1000 рублей в месяц или 6000 рублей в год, что для российского пользователя достаточно дорого. Приложение работает только на английском языке, инструкции к упражнениям не переведены. Интерфейс современный, но некоторым пользователям кажется перегруженным графиками и цифрами.

\subsubsection{JEFIT}

\noindent JEFIT - приложение с большим сообществом пользователей и огромной библиотекой упражнений. Оно существует с 2010 года и зарекомендовало себя как надёжный инструмент для ведения дневника тренировок. Основные характеристики:

\begin{itemize}
\item более 1300 упражнений с детальными анимациями и описаниями правильной техники;
\item возможность создавать свои собственные тренировки из имеющихся упражнений;
\item форум, где можно общаться с другими спортсменами, задавать вопросы и делиться опытом;
\item планировщик тренировок на дни и недели вперёд;
\item подробная статистика по каждому упражнению и мышечной группе.
\end{itemize}

\noindent \textbf{Недостатки:} устаревший дизайн, который не обновлялся последние несколько лет. Много рекламы в бесплатной версии - баннеры появляются после каждого подхода. Интерфейс на русском языке переведён не полностью, встречаются английские термины.

\medskip
\noindent \textbf{Вывод:} существующие решения либо платные (Fitbod, дорогая подписка MyFitnessPal), либо не учитывают индивидуальные особенности пользователя в полной мере (MyFitnessPal не даёт планов тренировок), либо имеют устаревший или перегруженный интерфейс (JEFIT, бесплатная версия MyFitnessPal). Проект FITGUIDE призван заполнить эту нишу: он будет полностью бесплатным, с современным дизайном и персонализированной генерацией планов на основе всех введённых параметров.

\subsection{Инструменты для разработки ML-интерфейсов}

\noindent Для создания клиентской части были рассмотрены несколько популярных инструментов. Критериями выбора были: простота освоения (чтобы студент второго курса мог разобраться за разумное время), возможности кастомизации дизайна (чтобы приложение не выглядело как стандартный шаблон), наличие готовых компонентов для работы с файлами и API.

\subsubsection{Flask}

\noindent Flask - это микро-фреймворк для веб-разработки на Python. Он очень лёгкий и гибкий, даёт разработчику полный контроль над приложением. С помощью Flask можно создать что угодно - от простой странички до сложного веб-сервиса.

\noindent Однако Flask требует знания HTML (языка разметки веб-страниц), CSS (каскадных таблиц стилей для оформления) и хотя бы базового JavaScript (языка программирования для браузеров). Для студента, который не изучал веб-технологии углублённо, Flask оказался слишком сложным. Чтобы сделать на Flask то же самое, что было сделано на выбранном фреймворке, потребовалось бы в десять раз больше времени и знаний.

\subsubsection{Django}

\noindent Django - это мощный и тяжеловесный фреймворк для создания крупных веб-проектов. Он включает в себя всё необходимое «из коробки»: ORM (систему для работы с базами данных без написания SQL-запросов), систему аутентификации пользователей (регистрация, вход, восстановление пароля), административную панель для управления сайтом.

\medskip
\noindent Однако для небольшого студенческого проекта Django - это «из пушки по воробьям». Он слишком тяжёлый (файлы проекта занимают десятки мегабайт), слишком сложный (кривая обучения очень крутая) и слишком медленный в разработке (даже простую форму нужно описывать в нескольких файлах). Для данной задачи Django категорически не подходит.
\vskip 5cm
\subsubsection{Gradio}

\noindent Gradio создан специально для быстрого прототипирования ML-приложений. Буквально в пять-десять строк кода можно получить работающий интерфейс с полями ввода и выводом результата. Gradio автоматически создаёт веб-страницу, настраивает сервер и обрабатывает запросы.

\medskip
\noindent Однако кастомизация дизайна в Gradio сильно ограничена. Все приложения на Gradio выглядят одинаково и напоминают стандартные блокноты Jupyter Notebook: белый фон, стандартные синие кнопки, нет возможности добавить свои анимации или изменить расположение элементов. Хотелось сделать что-то более современное и красивое, с тёмной темой и бирюзовыми акцентами. Gradio для этого не подходит.

\subsubsection{Streamlit}

\noindent Streamlit - это золотая середина между простотой Gradio и гибкостью Flask. С одной стороны, он так же прост, как Gradio: интерфейс описывается на чистом Python, не нужно писать HTML или JavaScript. С другой стороны, Streamlit позволяет добавлять свой CSS, анимации, настраивать расположение элементов с помощью колонок и контейнеров.

\medskip
\noindent Кроме того, у Streamlit большое сообщество и множество готовых компонентов. Например, для данной задачи очень пригодился компонент для скачивания файлов — достаточно одной строчки кода, и пользователь может сохранить сгенерированный план на компьютер.

\medskip
\noindent \textbf{Был выбран Streamlit.} В процессе работы это решение полностью оправдало себя. Streamlit оказался именно тем инструментом, который позволяет быстро разрабатывать и при этом получать красивый результат.
\vskip 5cm
\subsection{Выбор нейросетевой модели и API}

\noindent Для генерации фитнес-планов нужна языковая модель, которая соответствует следующим критериям:

\begin{enumerate}
\item хорошо понимает русский язык (большинство пользователей приложения - русскоязычные);
\item способна генерировать связные тексты объёмом от пятисот до тысячи символов (примерно один абзац текста);
\item доступна бесплатно для некоммерческого использования;
\item имеет удобный API с хорошей документацией.
\end{enumerate}

\noindent Были рассмотрены несколько вариантов, которые доступны на сегодняшний день.

\subsubsection{OpenAI GPT (ChatGPT)}

\noindent Самая известная языковая модель в мире. GPT-4 (последняя версия на момент написания работы) отлично понимает русский язык, генерирует очень качественные ответы, почти не ошибается в грамматике и стиле. Однако API OpenAI платный. Стоимость одного запроса объёмом около тысячи символов составляет примерно два-три рубля. Для студенческого проекта это пусть и небольшие, но всё же траты, которых хотелось избежать. Кроме того, требуется привязка банковской карты, что проблематично.

\subsubsection{YandexGPT}

\noindent Отечественная модель от компании Яндекс. У неё есть бесплатные тарифы для разработчиков, она хорошо знает русский язык (это её родной язык, в отличие от GPT, который обучался в основном на английском). Однако для доступа к API Яндекса нужно проходить довольно сложную процедуру регистрации: создавать платёжный аккаунт, подтверждать личность, разбираться в их облачной платформе Yandex Cloud. Для студенческого проекта это показалось излишне сложным.

\subsubsection{Hugging Face (модель Llama 3)}

\noindent Hugging Face - это платформа, на которой собраны десятки тысяч моделей машинного обучения. Любой разработчик может выложить свою модель, а другие разработчики могут использовать её через API. Модель под названием «Llama 3» была разработана компанией Meta (Facebook) и считается одной из лучших открытых языковых моделей.

\medskip
\noindent Преимущества этого варианта:
\begin{itemize}
\item полностью бесплатно для некоммерческого использования;
\item до пятисот бесплатных запросов в день - этого с запасом хватит для проекта;
\item хорошо понимает русский язык, если в запросе специально попросить отвечать по-русски;
\item API очень простой, хорошо документированный.
\end{itemize}

\medskip
\noindent \textbf{Был выбран Hugging Face и модель Llama 3.} Это решение оказалось оптимальным по соотношению цена/качество. Модель бесплатна, с ней легко работать, а качество генерации на русском языке после некоторой настройки промптов стало вполне приемлемым.

\subsection{Итоговый технологический стек}

\noindent Таким образом, для реализации проекта используется следующий набор технологий:

\begin{itemize}
\item \textbf{Язык программирования:} Python версии 3.10. Python выбран потому, что он прост в освоении, имеет огромное количество библиотек, и на нём удобно писать как логику приложения, так и взаимодействие с API.

\item \textbf{Фреймворк для интерфейса:} Streamlit версии 1.28 или новее. Streamlit позволяет быстро создавать веб-интерфейсы на чистом Python, поддерживает кастомизацию через CSS.

\item \textbf{API для доступа к нейросети:} Hugging Face Inference API. Это сервис, который предоставляет доступ к сотням моделей, включая Llama 3, через простой HTTP-интерфейс.

\item \textbf{Нейросетевая модель:} meta-llama/Meta-Llama-3-8B-Instruct. Это инструктивная версия Llama 3, оптимизированная для диалогов и выполнения инструкций.

\item \textbf{Дополнительные библиотеки:}
\begin{itemize}
\item python-dotenv - для безопасного хранения токена доступа
к API в отдельном файле;
\item base64 - для преобразования изображений в текстовый формат, чтобы встраивать их в CSS;
\item random - для случайного выбора примера цели из заранее подготовленного списка;
\item unicodedata - для очистки текста ответа от невидимых символов, которые иногда генерирует нейросеть.
\end{itemize}
\end{itemize}

\newpage
\section{ПРОЕКТИРОВАНИЕ КЛИЕНТСКОЙ ЧАСТИ FITGUIDE}

\subsection{Функциональные требования к системе}

\noindent Перед тем как писать код, были подробно описаны требования к приложению. Этот список функциональных требований служил дорожной картой на протяжении всей разработки. Требования разбиты на несколько категорий.

\subsubsection{Сбор физических параметров}

\noindent Приложение должно предоставлять пользователю возможность ввести следующие параметры:

\begin{itemize}
\item \textbf{Вес.} Поле для ввода числа от 40 до 200 килограммов с шагом 1 килограмм. Значение по умолчанию - 75 кг.

\item \textbf{Рост.} Поле для ввода числа от 100 до 250 сантиметров с шагом 1 сантиметр. Значение по умолчанию - 175 см.

\item \textbf{Пол.} Выбор из двух вариантов: «Мужской» или «Женский». Значение по умолчанию - «Мужской».

\item \textbf{Тип телосложения.} Выбор из трёх вариантов: «Эктоморф» (худощавое телосложение, трудно набирает массу), «Мезоморф» (атлетическое телосложение, легко набирает мышцы), «Эндоморф» (склонность к полноте, трудно худеть). Значение по умолчанию - «Мезоморф».

\item \textbf{Уровень опыта.} Ползунок с тремя позициями: «Новичок» (занимается менее полугода), «Средний» (от полугода до двух лет), «Профи» (более двух лет). Значение по умолчанию - «Средний».

\item \textbf{Доступный инвентарь.} Множественный выбор из четырёх вариантов: «Зал» (полноценный тренажёрный зал), «Гантели» (только гантели дома), «Турник» (турник и брусья), «Без инвентаря» (тренировки с собственным весом). Значение по умолчанию - «Без инвентаря».
\end{itemize}

\subsubsection{Ввод цели}

\noindent Приложение должно предоставлять текстовое поле для ввода цели свободным текстом. Поле должно иметь подсказку-пример, которая случайным образом выбирается из заранее подготовленного списка. Примеры целей: «набрать 5 кг мышц к лету», «подсушиться и убрать живот к отпуску», «похудеть на 8 кг за 3 месяца».

\subsubsection{Генерация плана}

\noindent При нажатии на кнопку «Сформировать фитнес-план» приложение должно:
\begin{enumerate}
\item Проверить, что поле цели не пустое. Если пустое - показать предупреждение.
\item Показать индикатор загрузки (спиннер) с текстом «Генерируем план...» для обратной связи с пользователем.
\item Сформировать промпт, включив в него все параметры пользователя.
\item Отправить запрос к Hugging Face API.
\item Получить ответ, очистить его от невидимых символов.
\item Отобразить ответ в специальном оформленном блоке.
\item Показать кнопку для сохранения плана в файл.
\end{enumerate}

\subsubsection{Сохранение плана}

\noindent После успешной генерации плана должна появляться кнопка «Сохранить план в TXT». При нажатии на неё браузер должен предложить пользователю сохранить текстовый файл с именем «my_fitplan.txt», содержащим сгенерированный план.
\vskip 5cm
\subsubsection{Совет дня}

\noindent При загрузке страницы (каждый раз, когда пользователь открывает или обновляет приложение) должен генерироваться случайный полезный совет по фитнесу. Совет должен быть коротким (одно-два предложения). Если API недоступно, должен показываться заранее заготовленный резервный совет.

\subsection{Разработка макетов интерфейса}

\noindent Чтобы чётко представлять, как будет выглядеть приложение, сначала был нарисован макет на бумаге, а затем перенесён в Figma - бесплатный онлайн-инструмент для создания макетов интерфейсов.

\subsubsection{Общая структура}

\noindent Экран разделён на две основные области:

\begin{enumerate}
\item \textbf{Боковая панель (сайдбар)} - расположена слева, занимает примерно 25% ширины экрана. Здесь находятся все параметры: вес, рост, пол, тип телосложения, опыт, инвентарь.

\item \textbf{Основная область} - расположена справа, занимает оставшиеся 75% ширины. Основная область, в свою очередь, разделена на три части:
\begin{itemize}
\item \textbf{Верхняя левая часть} - поле для ввода цели.
\item \textbf{Верхняя правая часть} - блок «Совет дня».
\item \textbf{Нижняя часть} - блок для отображения сгенерированного плана (изначально пустой, появляется после нажатия на кнопку).
\end{itemize}
\end{enumerate}

\noindent Такое расположение выбрано не случайно. Параметры пользователь вводит один раз за сессию, их можно «спрятать» в сайдбар, чтобы они не занимали лишнего места на экране. Цель же пользователь может менять часто (например, сначала хочу набрать массу, потом - похудеть), поэтому поле для её ввода находится на самом видном месте.

\subsubsection{Цветовая схема}

\noindent Цвета подбирались долго. В итоге выбрана бирюза на тёмном фоне - это выглядит дорого и современно. Плюс бирюза ассоциируется со здоровьем и энергией, для фитнес-приложения это самое то.

\begin{itemize}
\item \textbf{Основной цвет фона:} тёмно-зелёный с кодом #0F1C1A. Этот цвет создаёт ощущение глубины и спокойствия, не утомляет глаза при длительном использовании.

\item \textbf{Акцентный цвет (кнопки, рамки, важные элементы):} бирюзовый с кодом #5CE1D6. Этот цвет символизирует здоровье, энергию, движение вперёд.

\item \textbf{Цвет текста:} белый (#FFFFFF) и светло-серый (#CCCCCC) для хорошей читаемости на тёмном фоне.

\item \textbf{Цвет фона для блоков с планом и советом:} тёмно-зелёный с прозрачностью 92% (значение rgba(26,42,39,0.92)). Это создаёт эффект «матового стекла» — блоки как бы парят над фоном.
\end{itemize}

\subsubsection{Типографика}

\noindent Для заголовка приложения выбран крупный жирный шрифт. FITGUIDE написано белым цветом, а буквы AI - бирюзовым, чтобы подчеркнуть, что это приложение с искусственным интеллектом. Для основного текста оставлен стандартный шрифт Streamlit - он достаточно читаемый и хорошо работает на всех устройствах.

\subsubsection{Анимации}

\noindent Статичный интерфейс - скучно. Поэтому добавлены анимации. Логотип в сайдбаре пульсирует: мягко светится и затухает каждые три секунды. Так он привлекает внимание, но не раздражает.
\vskip 5cm
\subsection{Сценарии использования}

\noindent Для того чтобы лучше понять, как пользователи будут взаимодействовать с приложением, описаны несколько типичных сценариев использования. Это помогло проверить, все ли функции реализованы и не возникает ли где-то логических противоречий.

\subsubsection{Сценарий 1: Новый пользователь}

\noindent Пользователь впервые открывает приложение. Он видит боковую панель с параметрами, поле для цели и совет дня.

\begin{enumerate}
\item Пользователь заполняет параметры: вес 80 кг, рост 180 см, пол - мужской, телосложение - мезоморф, опыт - средний, инвентарь - зал.

\item В поле цели пишет: «Набрать 5 кг мышечной массы за 3 месяца».

\item Нажимает кнопку «СФОРМИРОВАТЬ ФИТНЕС-ПЛАН».

\item Появляется спиннер загрузки с текстом «Генерируем план...».

\item Через 3–4 секунды появляется блок с планом тренировок и питания, а также кнопка «Сохранить план».

\item Пользователь читает план, нажимает «Сохранить план» и сохраняет его на компьютер.
\end{enumerate}

\subsubsection{Сценарий 2: Опытный пользователь, меняющий цель}

\noindent Пользователь уже знаком с приложением, его параметры уже заполнены (если он не закрывал вкладку браузера).

\begin{enumerate}
\item Пользователь решает сменить цель с набора массы на сушку.
\item Стирает старую цель и пишет новую: «Подсушиться к лету, убрать живот».
\item Снова нажимает кнопку.
\item Получает новый план, соответствующий новой цели.
\end{enumerate}

\subsubsection{Сценарий 3: Ошибка API}

\noindent Пользователь пытается сгенерировать план, но API Hugging Face временно недоступно (например, превышен лимит запросов).

\begin{enumerate}
\item Пользователь заполняет параметры и нажимает кнопку.
\item Приложение пытается отправить запрос, но получает ошибку.
\item Вместо падения приложения пользователь видит понятное сообщение: «Ошибка при обращении к нейросети. Попробуйте позже или измените формулировку цели».
\item Блок «Совет дня» продолжает работать, показывая резервный совет, не зависящий от API.
\end{enumerate}

\newpage
\section{РЕАЛИЗАЦИЯ КЛИЕНТСКОЙ ЧАСТИ И ТЕСТИРОВАНИЕ}

\subsection{Организация проекта и управление состоянием}

\noindent Для удобства разработки и последующего сопровождения все файлы проекта организованы в одну директорию. Это упрощает запуск и распространение приложения. Главный исполняемый файл содержит весь код клиентской части. В отдельном файле хранится секретный токен доступа к API - этот файл не включается в систему контроля версий и не публикуется в открытый доступ.

\medskip
\noindent Ключевой технической проблемой, возникшей при разработке, было сохранение введённых пользователем данных между перезапусками скрипта. Как уже объяснялось в теоретической части, выбранный фреймворк перезапускает скрипт при каждом действии пользователя. Если не принять специальных мер, введённая цель будет теряться после нажатия на кнопку генерации.

\medskip
\noindent Для решения этой проблемы использован механизм хранения состояния сессии. При инициализации приложения проверяется, существует ли переменная для хранения цели в состоянии сессии. Если нет - создаётся с пустым значением. Затем поле ввода цели связывается с этой переменной. При каждом изменении текста в поле ввода переменная состояния обновляется. Благодаря этому, даже после перезапуска скрипта (который происходит при нажатии на кнопку), введённая цель остаётся на месте.

\subsection{Взаимодействие с API и обработка ответов}

\noindent Для подключения к нейросети через API используется официальная библиотека, предоставляемая Hugging Face. Токен доступа загружается из файла с переменными окружения - это безопаснее, чем хранить токен прямо в коде.

\medskip
\noindent При нажатии на кнопку генерации плана происходит следующее. Сначала проверяется, что поле цели не пустое. Если пользователь забыл ввести цель, появляется предупреждение. Если цель введена, формируется промпт - текстовая инструкция для нейросети. В промпте объясняется модели, кем она должна себя считать (элитный тренер FitGuide), какие данные о пользователе у неё есть (вес, рост, пол, тип телосложения, опыт, цель, инвентарь), и что от неё требуется (составить план тренировок и питания на русском языке).

\medskip
\noindent Затем промпт отправляется к API. Весь код, связанный с отправкой запроса, обёрнут в конструкцию перехвата ошибок. Это означает, что если что-то пойдёт не так (нет интернета, превышен лимит запросов, сервер Hugging Face временно недоступен), приложение не упадёт с ошибкой, а покажет пользователю понятное сообщение.

\medskip

\noindent Когда ответ от нейросети получен, он может содержать невидимые символы, которые мешают при копировании текста. Ответ очищается от этих символов с помощью специальной библиотеки, которая нормализует текст и заменяет проблемные символы на обычные пробелы.

\subsection{Кастомизация внешнего вида}

\noindent Стандартный интерфейс выбранного фреймворка имеет довольно спартанский, «технический» вид: серые тона, стандартные синие кнопки, отсутствие теней и анимаций. Чтобы приложение выглядело современно и привлекательно, добавлены собственные стили с помощью каскадных таблиц стилей (CSS).

\medskip
\noindent CSS встраивается непосредственно в код приложения через специальную функцию, которая позволяет вставлять HTML и CSS на страницу. Вот основные изменения, которые были сделаны:

\begin{itemize}
\item \textbf{Фон.} Вместо скучного белого или светло-серого фона установлен тёмно-зелёный градиент, поверх которого наложены три изображения, медленно сменяющие друг друга. Изображения закодированы в текстовый формат прямо в коде, чтобы не хранить отдельные файлы на сервере.

\item \textbf{Кнопки.} Стандартные кнопки перекрашены в бирюзовый цвет, который является акцентным цветом всего приложения. При наведении курсора кнопка становится ярче и начинает светиться.

\item \textbf{Контейнеры для плана и совета.} Блоки, в которых отображается сгенерированный план и совет дня, имеют полупрозрачный тёмный фон, бирюзовую рамку, скруглённые углы и лёгкую тень. Это создаёт эффект «всплывающей карточки».

\item \textbf{Логотип.} Логотип приложения пульсирует - его тень мягко расширяется и сужается с периодом три секунды, привлекая внимание пользователя.
\end{itemize}

\noindent Все эти изменения сделали приложение визуально привлекательным и запоминающимся.

\subsection{Генерация совета дня и обработка ошибок}

\noindent Совет дня генерируется при каждой загрузке страницы отдельным запросом к API. Запрос короткий - нейросеть просится дать один полезный совет по фитнесу максимум в два предложения.

\medskip
\noindent Если API недоступно или возникает ошибка, пользователь не видит красного сообщения на английском языке. Вместо этого блок совета дня показывает заранее заготовленный текст: «Регулярность важнее интенсивности. Тренируйтесь 3–4 раза в неделю». Это пример так называемой «изящной деградации» - приложение продолжает работать в урезанном режиме и остаётся полезным для пользователя.

\subsection{Тестирование и выявленные проблемы}

\noindent Тестирование проводилось вручную. Приложение проверялось на разных устройствах (ноутбук, планшет, телефон), в разных браузерах (Google Chrome, Yandex Browser) и при разных сценариях использования.

\subsubsection{Проблема с потерей цели}

\noindent \textbf{Описание:} при нажатии на кнопку «Сформировать план» введённая цель исчезала из поля ввода. Пользователю приходилось вводить её заново перед каждой генерацией, что очень неудобно.

\medskip
\noindent \textbf{Причина:} скрипт перезапускался при нажатии на кнопку, переменная, хранящая цель, создавалась заново с пустым значением.

\medskip
\noindent \textbf{Решение:} использование механизма хранения состояния сессии для цели. Теперь цель сохраняется между перезапусками.

\subsubsection{Проблема с невидимыми символами}

\noindent \textbf{Описание:} при копировании текста сгенерированного плана в блокнот или Word появлялись лишние пробелы и странные переносы строк.

\medskip
\noindent \textbf{Причина:} нейросеть иногда генерирует специальные невидимые символы, которые не отображаются на экране, но влияют на форматирование при копировании.

\medskip
\noindent \textbf{Решение:} очистка текста через нормализацию Юникода и замену проблемных символов на обычные пробелы.

\subsubsection{Проблема с английским языком в ответах}

\noindent \textbf{Описание:} несмотря на явную просьбу отвечать на русском языке, в некоторых случаях нейросеть выдавала ответ на английском.

\medskip
\noindent \textbf{Причина:} модель Llama 3 обучалась в основном на английских текстах, английский для неё «родной».

\medskip
\noindent \textbf{Решение:} усиление требования в промпте, добавление примера русскоязычного ответа. В большинстве случаев это помогло.
\vskip 5cm
\subsubsection{Проблема с дёрганьем фона}

\noindent \textbf{Описание:} при смене фонового изображения заметно «дёрганье» - изображение как будто передёргивается.

\medskip
\noindent \textbf{Причина:} слишком большое количество ключевых кадров в анимации и слишком быстрая смена изображений.

\medskip
\noindent \textbf{Решение:} уменьшено количество кадров с шести до трёх, увеличены временные интервалы между сменами изображений.

\subsection{Результаты тестирования}

\noindent После устранения всех выявленных проблем приложение работает стабильно. Среднее время ответа API составляет 2–4 секунды, что пользователи воспринимают как приемлемое. Интерфейс корректно отображается на всех протестированных устройствах. Функция сохранения плана создаёт файл с правильной кодировкой UTF-8, русские буквы отображаются нормально.

\newpage
\addcontentsline{toc}{section}{ЗАКЛЮЧЕНИЕ}
\section*{ЗАКЛЮЧЕНИЕ}

\noindent В ходе выполнения курсовой работы была спроектирована и реализована клиентская часть веб-приложения FITGUIDE с поддержкой интерактивных компонентов машинного обучения. Все поставленные задачи выполнены в полном объёме.

\subsection*{Основные результаты работы}

\noindent \textbf{Теоретические результаты:}
\begin{enumerate}
\item Изучены современные подходы к созданию пользовательских интерфейсов для ML-приложений. Выявлены ключевые особенности: необходимость индикации загрузки из-за неопределённого времени ответа, обработка ошибок API, персонализация ввода как главное преимущество.

\item Проведён сравнительный анализ существующих решений для фитнес-планирования (MyFitnessPal, Fitbod, JEFIT). Выявлены их недостатки: платность, отсутствие персонализации, устаревший дизайн.

\item Проведён сравнительный анализ инструментов быстрой разработки интерфейсов на Python (Flask, Django, Gradio, Streamlit). Обоснован выбор Streamlit как наиболее подходящего фреймворка для студенческого проекта.

\item Обоснован выбор Hugging Face API и модели Llama 3 как оптимального решения по критериям «бесплатно - просто - достаточно качественно».

\end{enumerate}

\noindent \textbf{Практические результаты:}
\begin{enumerate}
\item Разработана структура интерфейса, созданы макеты страниц. Продумана цветовая схема и анимации.

\item Реализован программный код клиентской части на языке Python с использованием фреймворка Streamlit.

\item Реализованы все функциональные требования: сбор шести параметров пользователя, ввод цели свободным текстом, генерация плана через API, сохранение в файл, совет дня.

\item Выполнена кастомизация внешнего вида через CSS: анимированный фон, бирюзовая цветовая гамма, эффекты при наведении, пульсирующий логотип.

\item Проведено тестирование, выявлены и исправлены четыре основные проблемы (потеря цели, невидимые символы, английский язык в ответах, дёрганье фона).
\end{enumerate}

\subsection*{Приобретённые навыки}

\noindent Работа над проектом дала ценный практический опыт, который пригодится в дальнейшем обучении и, возможно, в профессиональной деятельности:

\begin{itemize}
\item Получены навыки работы с API - отправка HTTP-запросов, обработка ответов, обработка ошибок.

\item Освоен фреймворк Streamlit - от простых виджетов (кнопки, поля ввода, ползунки) до сложной кастомизации через CSS.

\item Понята концепция состояния приложения и необходимость механизма хранения сессии для сохранения данных между перезапусками.

\item Освоены методы решения реальных проблем (невидимые символы, потеря состояния, адаптивность, обработка ошибок).

\item Приобретены навыки проектирования интерфейса с продумыванием сценариев использования и потребностей разных категорий пользователей.
\end{itemize}
\vskip 5cm
\subsection*{Направления дальнейшего развития}

\noindent Разработанное приложение является полностью рабочим прототипом, но его можно улучшить по следующим направлениям:

\begin{enumerate}
\item \textbf{Регистрация и профили пользователей.} Добавить возможность создавать аккаунты, чтобы каждый пользователь мог сохранять свою историю планов и прогресс.

\item \textbf{Редактирование плана.} Добавить возможность редактировать сгенерированный план прямо в интерфейсе, а не только сохранять.

\item \textbf{Визуализация прогресса.} Добавить графики изменения веса, объёмов, силы - это сделает приложение ещё полезнее.

\item \textbf{Локальная модель.} Настроить запуск нейросети на локальном компьютере пользователя, чтобы приложение могло работать без интернета.

\item \textbf{Мобильное приложение.} Написать нативное приложение для iOS и Android, так как веб-версия на телефоне работает, но не идеально.
\end{enumerate}

\subsection*{Заключительная оценка}

\noindent Разработанное приложение FITGUIDE является полностью работоспособным, удобным и визуально привлекательным. Оно решает поставленную задачу - помогает пользователям получать персонализированные
фитнес-планы с учётом их индивидуальных параметров и цели. Проект может быть использован как прототип для дальнейшего коммерческого развития, а также как учебный пример для студентов, изучающих машинное обучение и веб-разработку.

\medskip

\newpage
\renewcommand\refname{\centering СПИСОК ИСПОЛЬЗОВАННЫХ ИСТОЧНИКОВ}
\clearpage
\addcontentsline{toc}{section}{СПИСОК ИСПОЛЬЗОВАННЫХ ИСТОЧНИКОВ}
\begin{thebibliography}{99}

\bibitem{streamlit_docs} Streamlit Documentation. Основы работы с session
state и кастомизация. [Электронный ресурс] URL: https://docs.streamlit.io/
(дата обращения: 10.04.2026).

\bibitem{hf_docs} Hugging Face Hub Docs. Inference API: быстрое
прототипирование. [Электронный ресурс] URL:
https://huggingface.co/docs/api-inference/index (дата обращения: 12.04.2026).

\bibitem{chertkova} Черткова Е. А. Программная инженерия. Визуальное
моделирование программных систем: учебник для вузов. - 2-е изд. -
Москва: Юрайт, 2022. — 147 с.

\bibitem{zaramenskih} Зараменских Е. П. Разработка информационных систем:
учебник и практикум. - 2-е изд. - Москва: Юрайт, 2025. - 78 с.

\bibitem{galias} Галиаскаров Э. Г. Анализ и проектирование систем с
использованием UML: учебное пособие. - Москва: Юрайт, 2022. - 125 с.

\bibitem{ratush} Ратушняк Г. Я., Золкин А. Л. Технологии разработки и
проектирования информационных систем. Часть 1. - Москва: Русайнс,

201 с.

\bibitem{vereschagina} Верещагина Е. А., Золкин А. Л. Методы и алгоритмы
организации процессов мониторинга информационных систем. - Москва:
РУСАЙНС, 2024. - 152 с.

\bibitem{isaev} Исаев Г. Н. Теоретико-методологические основы качества
информационных систем. - Москва: ИНФРА-М, 2024. - 294 с.

\bibitem{habr} Создаём чат-бота на Streamlit и Hugging Face за 15 минут.
Habr, 2024. [Электронный ресурс] URL: https://habr.com/ (дата обращения:
15.04.2026).

\bibitem{method} Методические рекомендации по подготовке курсовых работ
по направлению «Прикладная информатика». Липецк: Липецкий филиал
Финансового университета, 2024. - 35 с.

\bibitem{lipetsk} Учебно-методическое пособие по выполнению курсовых
проектов для студентов направления 09.03.03. - Липецк: Издательство
Липецкого филиала Финуниверситета, 2024. - 48 с.

\end{document}

